\documentclass[11pt]{article}
\usepackage[margin=1in]{geometry}
\usepackage{amsmath,amssymb,amsthm,hyperref}
\begin{document}
\begin{center}{\Large Degree-zero rational cohomology does not decay}\\[4pt]
Disproof of conjecture 00000001646\\[4pt]
ziangni-sys --- October 8, 2026\end{center}

The conjecture asserts that
$\dim_{\mathbb Q}H^i(\operatorname{Mod}(\Sigma_g);\mathbb Q)$ tends to zero
exponentially as $g$ grows, whenever $i\le (2g-2)/3$. Both language versions
include degree $i=0$: no positive-degree or reduced-cohomology restriction
is stated. Ordinary group cohomology in this degree gives a counterexample.

\paragraph{Degree-zero calculation for every group.}
Let $G$ be any group and give $\mathbb Q$ the trivial $G$-action. In the
inhomogeneous cochain complex,
\[
 C^0(G;\mathbb Q)=\mathbb Q,\qquad
 C^1(G;\mathbb Q)=\{f:G\to\mathbb Q\},\qquad
 (d^0q)(h)=h\cdot q-q.
\]
Since the action is trivial, $d^0=0$. Thus every rational number is a
zero-cocycle. There are no negative-degree cochains, so the zero-coboundary
space is $\operatorname{im}(d^{-1})=0$. The actual cohomology quotient is
therefore
\[
 H^0(G;\mathbb Q)=\ker(d^0)/\operatorname{im}(d^{-1})
 =\mathbb Q/0\cong\mathbb Q,
 \qquad \dim_{\mathbb Q}H^0(G;\mathbb Q)=1.
\]
This calculation applies to every mapping class group, irrespective of genus,
connectedness conventions for surfaces or the choice between standard
mapping class group variants.

\paragraph{Failure in the claimed stable range.}
For every integer $g\ge1$ we have $0\le(2g-2)/3$. Hence $i=0$ belongs to the
conjecture's stable range at all sufficiently large genera. Nevertheless,
\[
 \dim_{\mathbb Q}H^0(\operatorname{Mod}(\Sigma_g);\mathbb Q)=1
 \quad\text{for every }g.
\]
The sequence has limit one and cannot have limit zero, much less exponential
decay with constant $-\log2$. Equivalently, the neighborhood
$(-\tfrac12,\tfrac12)$ of zero contains none of its terms.

\paragraph{Scope.}
The introductory sentence mentions $l^2$-Betti numbers, whereas the claimed
quantitative formula explicitly uses ordinary rational cohomology
$H^i(-;\mathbb Q)$. These are different invariants. The counterexample
refutes the displayed ordinary-cohomology claim as written. It neither
asserts nor refutes any general $l^2$-vanishing theorem. Replacing $H^0$
by reduced cohomology or explicitly restricting to positive degrees would
change the statement and would require a separate analysis.

\paragraph{Lean certificate.}
The project constructs trivial-action cochains, the degree-zero
differential, its cocycle kernel, negative-degree zero cochains and the
boundary image. An explicit linear equivalence identifies the quotient
with $\mathbb Q$, proving dimension one for every group. It verifies
stable-range inclusion and proves nonconvergence to zero for every group
family, including when dimensions are only specified on the stable tail.
\end{document}
